\section*{Chapitre \ref{ch:b2:energy-geology} --- Géologie énergétique}
% numbers from the Book 4 second-writer scratch calculator calc/ch30.py

\begin{solution}{exo:b2:energy-geology:1}
(a) 1.58~: type~I~; 0.95~: type~III~; 1.24~: type~II~; 1.12~: mélange de types~II/III. Le kérogène de type~III, à
0.95, produit surtout du gaz. (b) H/C $= 150/100 = 1.50$, O/C $= 8/100 = 0.08$. (c) Pour \qty{100}{g}~:
$77.4/12.011 = \qty{6.444}{mol}$ de C, $9.8/1.008 = \qty{9.722}{mol}$ de H, $8.2/15.999 =
\qty{0.513}{mol}$ d'O~; H/C $= 1.51$ et O/C $= 0.0795$, en accord avec la formule.
\end{solution}

\begin{solution}{exo:b2:energy-geology:2}
(a) 0.45~: immature, en dessous du seuil à \qty{0.60}{\percent}~; 0.68~: génération d'huile~; 1.1~: encore dans la fenêtre à huile, près de sa fin~;
1.6~: au-delà de 1.20 à \qty{1.30}{\percent}, l'huile se décompose en méthane et la roche produit du gaz~;
2.5~: gaz seul, la génération se poursuivant jusqu'à environ \qty{4.8}{\percent}. (b) Non~: la génération de méthane à partir du type~III commence vers \qty{0.73}{\percent}.
(c) $F = (\ln R_o + 1.6)/3.7$~: $(-0.511 +
1.6)/3.7 = 0.294$ et $(0.262 + 1.6)/3.7 = 0.503$.
\end{solution}

\begin{solution}{exo:b2:energy-geology:3}
(a) Dans le réservoir, $6 \times 10^6 \times 25 \times 0.18 \times 0.70 = \qty{1.89e7}{m^3}$~;
$\mathrm{STOIIP} = 1.89 \times 10^7/1.25 = \qty{1.51e7}{m^3}$. (b) $0.30 \times 1.512 \times 10^7 =
\qty{4.54e6}{m^3}$ et $0.40 \times 1.512 \times 10^7 = \qty{6.05e6}{m^3}$.
\end{solution}

\begin{solution}{exo:b2:energy-geology:4}
(a) $V = 10^7 \times 200 = \qty{2e9}{m^3}$ et $q_R = 2.5 \times 10^6 \times 2 \times 10^9 \times 60
= \qty{3.0e17}{J}$. (b) En dessous de \qty{90}{\celsius}~: une \omterm{def:b2:energy-geology:geothermal}{ressource basse énergie}, pour l'usage direct de la
chaleur. (c) $P = 0.25 \times 3.0 \times 10^{17}/(30 \times \qty{3.156e7}{s}) = \qty{79.2}{MW}$ de puissance thermique.
\end{solution}

\begin{solution}{exo:b2:energy-geology:5}
(a) $k_2/k_1 = \exp[(E/R)(1/T_1 - 1/T_2)] = 2$ avec $1/373.15 - 1/383.15 = \qty{6.994e-5}{K^{-1}}$
donne $E = 8.314 \times \ln 2/(6.994 \times 10^{-5}) = \qty{82.4}{kJ/mol}$. (b) $\exp(26800 \times
6.994 \times 10^{-5}) = 6.52$. (c) La règle du doublement, dont l'énergie effective est plus faible, accorde
relativement plus de poids aux longues durées passées à basse température~;
pour un kérogène de type~I de forte énergie d'activation, la TTI annonce une génération trop précoce.
\end{solution}

\begin{solution}{exo:b2:energy-geology:6}
(a) $x = 1/2$ dans l'approximation donne $(ART^2/\beta E)\,e^{-E/RT} = \ln 2$~; en prenant les logarithmes,
$E/RT = \ln[AT^2/(\beta(E/R)\ln 2)]$, et la forme énoncée en découle. En partant d'une estimation voisine de
\qty{400}{K} et en itérant avec $\beta$ en kelvins par seconde~: à \qty{1}{\celsius\per\Ma}
($\beta = \qty{3.169e-14}{K/s}$), $T = \qty{409.6}{K}$, soit \qty{136.4}{\celsius}~; à \qty{10}{K} par
minute ($\beta = \qty{0.1667}{K/s}$, $\num{5.26e12}$ fois plus rapide), $T = \qty{719.0}{K}$,
soit \qty{445.9}{\celsius}. (b) La différence vaut \qty{309.5}{K}~; l'estimation $6.259 \times
\ln(5.26 \times 10^{12}) = 6.259 \times 29.29 = \qty{183}{K}$ est beaucoup trop petite, car $RT^2/E$
croît lui-même avec $T$ (jusqu'à \qty{19.3}{K} à \qty{719}{K})~: le décalage linéaire ne vaut que pour de petites
variations de vitesse.
\end{solution}

\begin{solution}{exo:b2:energy-geology:7}
$\lambda_5 - \lambda_8 = \qty{8.29725e-10}{an^{-1}}$. Avec $r_0 = 0.7204/99.2742 = 0.0072567$.
(a) À \qty{2.0}{\Ga}, $r = 0.0072567\,e^{1.6595} = 0.03814$, part $0.03814/1.03814 =
\qty{3.67}{\percent}$~; à \qty{1.8}{\Ga}, $r = 0.03231$, part \qty{3.13}{\percent}. (b) $r =
3/97 = 0.030928$ et $t = \ln(0.030928/0.0072567)/(8.29725 \times 10^{-10}) = \qty{1.75}{\Ga}$.
(c) Le minerai appauvri et l'uranium normal ont décru depuis par les mêmes facteurs, si bien que le rapport de leurs parts
en $^{235}$U est le même aujourd'hui que juste après les réactions~:
le minerai avait alors la moitié de la part normale, donc environ la moitié de son $^{235}$U avait été fissionnée.
\end{solution}

\begin{solution}{exo:b2:energy-geology:8}
(a) $10^9/1.872 = \qty{5.34e8}{m^3}$~; $10^9/621.4 = \qty{1.61e6}{m^3}$~; $10^9/686.4 =
\qty{1.46e6}{m^3}$. (b) $7 \times 10^9/583 = \qty{1.20e7}{m^3}$ de \ce{CO2}, qui remplit
$0.35 \times 0.30 = 0.105$ du volume de roche~: $1.20 \times 10^7/0.105 = \qty{1.14e8}{m^3}$ de roche,
soit une épaisseur moyenne de $1.143 \times 10^8/(5 \times 10^6) = \qty{22.9}{m}$.
\end{solution}

\begin{solution}{exo:b2:energy-geology:9}
(a) $V = \qty{9e9}{m^3}$ et $q_R = 2.6 \times 10^6 \times 9 \times 10^9 \times 235 =
\qty{5.50e18}{J}$. (b) Sur $30 \times \qty{3.156e7}{s} = \qty{9.47e8}{s}$~: $R_g = 0.1$ donne
\qty{581}{MW} et $0.25$ donne \qty{1450}{MW} de chaleur. (c) $6 \times (10 \text{ à } 20) = 60$ à
\qty{120}{MW} d'électricité, $60/581 = 0.10$ à $120/581 = 0.21$ de la puissance thermique~: l'essentiel de la
chaleur est rejeté, comme dans toute centrale travaillant avec un écart de température modeste.
\end{solution}

\begin{solution}{exo:b2:energy-geology:10}
(a) La résistance vaut $0.6 \times (80 - 39) = \qty{24.6}{MPa} > \qty{22}{MPa}$~: stable. La rupture exige
$0.6\,\Delta p = 24.6 - 22$, soit $\Delta p = \qty{4.33}{MPa} = \qty{43}{bar}$, du même ordre que les
\qty{32}{bar} qui bornaient les séismes de Denver. (b) $(24.6 - 24)/0.6 = \qty{1.0}{MPa} =
\qty{10}{bar}$~: une faille déjà proche de la rupture n'a besoin que d'une faible hausse. (c) Le fluide injecté élève la
pression dans une zone qui s'étend par diffusion à travers les roches perméables et les zones de failles, à des kilomètres
de distance au bout de plusieurs années~; là où elle rencontre une faille déjà proche de la rupture, une faible hausse suffit.
\end{solution}

\begin{solution}{exo:b2:energy-geology:11}
(a) Poser $x = e^{-u}$ avec $u = (ART^2/\beta E)\,e^{-E/RT}$. Un taux de transformation de
\qty{10}{\percent} correspond à $x = 0.9$, $u = \ln(10/9)$~; \qty{90}{\percent} à $u = \ln 10$. Comme
$\mathrm{d}\ln u/\mathrm{d}T = E/RT^2 + 2/T \approx E/RT^2$, $\Delta(\ln u) \approx (E/RT^2)\,\Delta T$,
ce qui donne la formule, avec $\ln\ln 10 - \ln\ln(10/9) = 0.834 + 2.250 = 3.084$. (b) Avec $T =
\qty{409.6}{K}$, $RT^2/E = \qty{6.26}{K}$ et $\Delta T = \qty{19.3}{K}$~;
la courbe exacte à énergie unique donne \qty{18.6}{K}, de 125.3 à \qty{143.8}{\celsius}. (c) La courbe pleine passe de
\qty{10}{\percent} vers \qty{122}{\celsius} à \qty{90}{\percent} vers \qty{148}{\celsius}, sur environ
\qty{26}{K}~: l'étalement des énergies d'activation élargit l'intervalle de conversion.
\end{solution}

\begin{solution}{exo:b2:energy-geology:12}
(a) $2 \times 10^{10}/20 = \qty{1.0e9}{m^3}$ (\qty{1.0}{km^3}) et $2 \times 10^{10}/60 =
\qty{3.3e8}{m^3}$. (b) $2 \times 10^{10}/621.4 = \qty{3.22e7}{m^3}$. (c) $3.22 \times 10^7/0.02 =
\qty{1.61e9}{m^3}$ de pores~; avec $50 \times 0.30 = \qty{15}{m}$ d'espace poreux par mètre carré, il faut
$1.61 \times 10^9/15 = \qty{1.07e8}{m^2}$, soit environ \qty{107}{km^2}.
\end{solution}

\begin{solution}{pb:b2:energy-geology:1}
\noindent\textbf{Partie I.}
\textbf{1.} À \qty{70}{\Ma}~: $50 \times 60 = \qty{3000}{m}$, $15 + 25 \times 3.0 =
\qty{90}{\celsius}$. Aujourd'hui~: $3000 + 70 \times 30 = \qty{5100}{m}$, \qty{142.5}{\celsius}.
\textbf{2.} $60 \times 0.025 = \qty{1.5}{\celsius\per\Ma}$, puis $30 \times 0.025 =
\qty{0.75}{\celsius\per\Ma}$.
\textbf{3.} De 15 à \qty{20}{\celsius}~: \qty{3.33}{\Ma}~; chaque intervalle de 20 à \qty{90}{\celsius}~:
$10/1.5 = \qty{6.67}{\Ma}$~; chaque intervalle de 90 à \qty{140}{\celsius}~: $10/0.75 = \qty{13.3}{\Ma}$~;
de 140 à \qty{142.5}{\celsius}~: \qty{3.33}{\Ma}.
\textbf{4.} Valeurs des intervalles $\Delta t_n 2^n$, entre parenthèses les sommes cumulées~: 15--20 ($n = -9$) 0.0065~; 20--30
0.026 (0.033)~; 30--40 0.052 (0.085)~; 40--50 0.104 (0.189)~; 50--60 0.208 (0.397)~; 60--70 0.417
(0.814)~; 70--80 0.833 (1.65)~; 80--90 1.67 (3.31)~; 90--100 6.67 (9.98)~; 100--110 13.3 (23.3)~;
110--120 26.7 (50.0)~; 120--130 53.3 (103)~; 130--140 107 (210)~; 140--142.5 ($n = 4$) 53.3~: TTI
aujourd'hui 263.
\textbf{5.} 15 est franchi dans l'intervalle 100--\qty{110}{\celsius}, à $100 + 10 \times
(15 - 9.98)/13.33 = \qty{103.8}{\celsius}$, atteint $(103.8 - 90)/0.75 = \qty{18.4}{\Ma}$ après
\qty{70}{\Ma}~: à \qty{51.6}{\Ma} et à $(103.8 - 15)/0.025 = \qty{3550}{m}$.
\textbf{6.} 75~: \qty{124.7}{\celsius}, \qty{23.7}{\Ma}, \qty{4390}{m}. 160~: \qty{135.3}{\celsius},
\qty{9.6}{\Ma}, \qty{4810}{m}.
\textbf{7.} Non~: à 263 elle a dépassé la fin de la génération d'huile (160), depuis environ \qty{10}{\Ma}.

\noindent\textbf{Partie II.}
\textbf{8.} $k = 9.21 \times 10^{13}\,e^{-26800/363.15} = \qty{8.20e-19}{s^{-1}}$. La roche est restée
en dessous de \qty{90}{\celsius} pendant les \qty{50}{\Ma} précédant \qty{70}{\Ma}, donc la fraction convertie est
inférieure à $k\,t = 8.20 \times 10^{-19} \times 50 \times 3.156 \times 10^{13} = \num{1.29e-3}$.
\textbf{9.} $\beta = 0.75/(3.156 \times 10^{13}) = \qty{2.376e-14}{K/s}$, $T = \qty{415.65}{K}$~:
$ART^2/(\beta E) = \num{2.498e28}$ et $e^{-E/RT} = \num{9.95e-29}$, donc $u = 2.486$, $x = 0.0833$ et
le taux de transformation vaut 0.917.
\textbf{10.} \qty{134.7}{\celsius}.
\textbf{11.} $70 - (134.7 - 90)/0.75 = \qty{10.4}{\Ma}$.
\textbf{12.} La TTI place le maximum de génération à \qty{23.7}{\Ma}, la cinétique \qty{13}{\Ma} plus tard,
parce que la règle du doublement correspond à $E \approx \qty{82}{kJ/mol}$ et pondère trop les basses températures
pour un kérogène de \qty{223}{kJ/mol} (exercice~\ref{exo:b2:energy-geology:5}).
\textbf{13.} Plus tôt~: une énergie plus faible abaisse la température de demi-conversion.

\noindent\textbf{Partie III.}
\textbf{14.} Roche~: $4 \times 10^8 \times 50 \times 2400 = \qty{4.8e13}{kg}$~; potentiel $0.030
\times 4.8 \times 10^{13} = \qty{1.44e12}{kg}$~; généré $0.92 \times 1.44 \times 10^{12} =
\qty{1.32e12}{kg}$.
\textbf{15.} $1.3248 \times 10^{12}/850 = \qty{1.56e9}{m^3}$.
\textbf{16.} $1.2 \times 10^7 \times 40 \times 0.22 \times 0.75/1.30 = \qty{6.09e7}{m^3}$.
\textbf{17.} $6.09 \times 10^7/1.56 \times 10^9 = 0.039$, \qty{3.9}{\percent}~: l'essentiel de l'huile générée
reste dans la roche mère, se perd le long des trajets de migration ou s'échappe vers la surface.
\textbf{18.} $0.35 \times 6.09 \times 10^7 = \qty{2.13e7}{m^3}$.
\textbf{19.} Environ les deux tiers de l'huile restent dans le réservoir après la récupération primaire et secondaire~;
l'injection de \ce{CO2} pourrait récupérer 10 à \qty{15}{\percent} de l'huile en place en plus.

\noindent\textbf{Partie IV.}
\textbf{20.} $15 + 25 \times 2 = \qty{65}{\celsius}$ et $0.101 + 1000 \times 9.807 \times 2000/10^6
= \qty{19.7}{MPa}$~: tous deux au-dessus du point critique (\qty{31.0}{\celsius}, \qty{7.38}{MPa}), si bien que le
\ce{CO2} est supercritique, à environ \qty{686}{kg.m^{-3}}.
\textbf{21.} $2.13 \times 10^7 \times 1.30 = \qty{2.77e7}{m^3}$.
\textbf{22.} Volume poreux $1.2 \times 10^7 \times 40 \times 0.22 = \qty{1.056e8}{m^3}$, et $E =
2.77 \times 10^7/1.056 \times 10^8 = 0.262$, bien au-dessus des 1 à \qty{4}{\percent} des formations salines~:
l'huile produite a vidé un espace connu et scellé.
\textbf{23.} $6 \times 10^8 \times 80 \times 0.20 \times 691 \times 0.02 = \qty{1.33e11}{kg}$,
\qty{133}{Mt}.
\textbf{24.} À \qty{1}{Mt} par an, environ 133 ans.
\textbf{25.} $M = \rho\,N_p\,B_o = 686 \times 2.77 \times 10^7 = \qty{1.90e10}{kg}$~: le réservoir épuisé peut stocker
\textbf{\qty[text-series-to-math=true]{19.0}{Mt}} de \ce{CO2}.
\end{solution}
